\chapter{Sufficiency of Certificates}
\label{ch:sufficiency}

Chapter~\ref{ch:proof-systems} showed that a task is served by a representation exactly when it is constant on the fibres of the extraction map (Proposition~\ref{prop:factor}). That statement is true of any function, and its content in the present setting depends entirely on which tasks are in play. This chapter makes the dependence precise. A \emph{task family} is declared, sufficiency is defined relative to it, and the resulting structure is shown to be a Galois connection on the partition lattice of the set of histories. Three consequences follow: a minimal sufficient representation exists and is unique up to isomorphism; sufficiency has a quantitative lower bound in cardinality and in entropy; and no representation that identifies two distinct histories is sufficient for every task. The last statement is the formal core of the claim that a verified result cannot be both maximally compressed and maximally reusable.

\section{Histories, representations and task families}

\begin{definition}[Setting]
\label{def:setting}
A \emph{history space} is a nonempty set $H$. A \emph{representation} is a function $\sigma:H\to R$. A \emph{task} with codomain $Y$ is a function $f:H\to Y$. A \emph{task family} is a set $\mathcal F$ of tasks, with codomains allowed to vary. The \emph{kernel} of a function $g$ on $H$ is the equivalence relation $\ker g=\{(h,h')\in H^2: g(h)=g(h')\}$.
\end{definition}

In the application, $H$ is the set of complete runs of a search procedure on a goal (the set $\mathrm{Run}(\varphi)$ of Chapter~\ref{ch:proof-systems}), and the canonical representation is the extraction map $\varrho:\mathrm{Run}(\varphi)\to D$ to the checked derivation. The abstraction is nevertheless useful because the same definitions apply to the obligation ledgers of Chapter~\ref{ch:obligations}, to proof terms modulo definitional equality (Chapter~\ref{ch:kernels}), and to normal forms (Chapter~\ref{ch:rewriting}).

\begin{definition}[Sufficiency]
\label{def:suff}
A representation $\sigma:H\to R$ is \emph{$f$-sufficient} if there is $f':R\to Y$ with $f=f'\circ\sigma$, and \emph{$\mathcal F$-sufficient} if it is $f$-sufficient for every $f\in\mathcal F$.
\end{definition}

By Proposition~\ref{prop:factor}, $\sigma$ is $f$-sufficient if and only if $\ker\sigma\subseteq\ker f$. The family enters through a single relation.

\begin{definition}[Family kernel]
\label{def:famker}
For a nonempty family $\mathcal F$ put $\ker\mathcal F=\bigcap_{f\in\mathcal F}\ker f$, that is, $h\sim_{\mathcal F}h'$ if and only if $f(h)=f(h')$ for all $f\in\mathcal F$. For the empty family $\ker\emptyset=H^2$.
\end{definition}

Equivalently, $\ker\mathcal F$ is the kernel of the joint task $\langle\mathcal F\rangle:H\to\prod_{f\in\mathcal F}Y_f$, $h\mapsto(f(h))_f$.

\begin{theorem}[Sufficiency is refinement]
\label{thm:suff-refine}
A representation $\sigma$ is $\mathcal F$-sufficient if and only if $\ker\sigma\subseteq\ker\mathcal F$. Consequently:
\begin{enumerate}[label=(\roman*),nosep]
\item the quotient map $q_{\mathcal F}:H\to H/\ker\mathcal F$ is $\mathcal F$-sufficient;
\item every $\mathcal F$-sufficient $\sigma$ factors $q_{\mathcal F}$, that is, $q_{\mathcal F}=u\circ\sigma$ for a unique $u:\sigma(H)\to H/\ker\mathcal F$; and
\item if $\mathcal F\subseteq\mathcal F'$ then every $\mathcal F'$-sufficient representation is $\mathcal F$-sufficient.
\end{enumerate}
\end{theorem}

\begin{proof}
If $\sigma$ is $\mathcal F$-sufficient then $\ker\sigma\subseteq\ker f$ for every $f\in\mathcal F$, hence $\ker\sigma\subseteq\bigcap_f\ker f$. Conversely, if $\ker\sigma\subseteq\ker\mathcal F\subseteq\ker f$, Proposition~\ref{prop:factor} gives sufficiency for $f$. For (i), $\ker q_{\mathcal F}=\ker\mathcal F$. For (ii), define $u(\sigma(h))=[h]_{\mathcal F}$; this is well defined because $\sigma(h)=\sigma(h')$ implies $(h,h')\in\ker\sigma\subseteq\ker\mathcal F$, and it is unique because $\sigma$ is onto $\sigma(H)$. For (iii), $\ker\mathcal F'\subseteq\ker\mathcal F$.
\end{proof}

Theorem~\ref{thm:suff-refine}(ii) says that $q_{\mathcal F}$ is the \emph{minimal sufficient representation}: it retains exactly the distinctions the family can detect and no others, and every sufficient representation retains at least these. The set of $\mathcal F$-sufficient kernels is the down-set of the partition lattice $\Pi(H)$ generated by $\ker\mathcal F$, ordered by refinement. In particular it is closed under common refinement (meet in the refinement order) and under passage to finer partitions, and it is \emph{not} closed under coarsening beyond $\ker\mathcal F$.

\begin{example}[A derivation as a representation]
\label{ex:deriv-rep}
Let $H$ consist of runs $h=(d,A)$ where $d$ is a checked derivation of a fixed goal and $A\subseteq\{1,\dots,k\}$ is the set of indices of alternative lemmas that were tried and refuted before $d$ was found. Take $\sigma(d,A)=d$. For the family $\mathcal F_{\mathrm{chk}}=\{h\mapsto\Chk(d,\varphi)\}$ the kernel $\ker\mathcal F_{\mathrm{chk}}$ identifies all runs with the same $d$, so $\sigma$ is sufficient and indeed minimal. For the family $\mathcal F_{\mathrm{ref}}=\{h\mapsto[i\in A]:1\le i\le k\}$ the kernel is $\{((d,A),(d',A')): A=A'\}$, which $\sigma$ does not refine whenever there are two runs with equal $d$ and different $A$. Each of the $2^k$ subsets $A$ is a separate class of $\ker\mathcal F_{\mathrm{ref}}$, so a sufficient representation for $\mathcal F_{\mathrm{ref}}$ has at least $2^k$ values (Theorem~\ref{thm:lower}).
\end{example}

\section{The Galois connection and the scope of a representation}

Declaring a family is one half of a duality; the other half asks what a given representation serves.

\begin{definition}[Served tasks and sufficient representations]
\label{def:galois}
For a representation $\sigma$ let $\mathrm{Serv}(\sigma)$ be the set of all tasks $f:H\to Y$ (any codomain) with $\ker\sigma\subseteq\ker f$. For a family $\mathcal F$ let $\mathrm{Rep}(\mathcal F)$ be the set of $\mathcal F$-sufficient representations.
\end{definition}

\begin{theorem}[Closure of a task family]
\label{thm:closure}
The pair $(\mathrm{Rep},\mathrm{Serv})$ is a Galois connection in the sense that $\mathcal F\subseteq\mathrm{Serv}(\sigma)$ if and only if $\sigma\in\mathrm{Rep}(\mathcal F)$. Moreover
\[
\mathrm{Serv}(\mathrm{Rep}(\mathcal F)):=\bigcap_{\sigma\in\mathrm{Rep}(\mathcal F)}\mathrm{Serv}(\sigma)=\{f:\ker\mathcal F\subseteq\ker f\},
\]
which is the set of tasks of the form $g\circ\langle\mathcal F\rangle$ for some function $g$. In words: a representation sufficient for $\mathcal F$ is automatically sufficient for every task computable from the joint values of $\mathcal F$, and for no other task in general.
\end{theorem}

\begin{proof}
The first statement is the definition of both sets. For the second, if $\ker\mathcal F\subseteq\ker f$ and $\sigma\in\mathrm{Rep}(\mathcal F)$ then $\ker\sigma\subseteq\ker\mathcal F\subseteq\ker f$ by Theorem~\ref{thm:suff-refine}, so $f\in\mathrm{Serv}(\sigma)$. Conversely, the quotient $q_{\mathcal F}$ lies in $\mathrm{Rep}(\mathcal F)$, so any $f$ in the intersection satisfies $\ker\mathcal F=\ker q_{\mathcal F}\subseteq\ker f$. Finally $\ker\mathcal F\subseteq\ker f$ holds if and only if $f=g\circ\langle\mathcal F\rangle$ for some $g$, by Proposition~\ref{prop:factor} applied to $\langle\mathcal F\rangle$.
\end{proof}

The closure operator $\mathcal F\mapsto\mathrm{Serv}(\mathrm{Rep}(\mathcal F))$ is idempotent, monotone and extensive, as for any Galois connection. Its practical content is a precise statement of what is gained and not gained by declaring a family. Declaring the checking task serves also every function of the checking verdict; declaring the checking task and the lemma-refutation tasks together serves all functions of the pair; but nothing declared serves a task whose answer depends on a distinction that none of the declared tasks detects. That a task fails to be in the closure is a decidable property when $H$ is finite and the tasks are given, and in general it is as hard as deciding the kernel inclusion.

\begin{remark}[Verified by enumeration]
For all history spaces with at most six elements, and for randomly drawn families of up to three tasks, the equivalence in Theorem~\ref{thm:suff-refine} and the characterization of the closure in Theorem~\ref{thm:closure} were confirmed by exhaustive enumeration of every representation up to relabelling (about three thousand families, no failures). This is a sanity check of the statements, not a proof; the proofs above are complete.
\end{remark}

\subsection*{Extending a representation to a larger family}

A result is certified once, against a family declared at the time. Later uses enlarge the family, and the question is what the cheapest extension is.

\begin{proposition}[Cheapest extension]
\label{prop:extend}
Let $\sigma$ be $\mathcal F$-sufficient and let $\mathcal F'$ be a further family. Among representations $\tau$ that are $\mathcal F'$-sufficient and that serve every task served by $\sigma$ (that is, $\mathrm{Serv}(\sigma)\subseteq\mathrm{Serv}(\tau)$), the coarsest has kernel $\ker\sigma\cap\ker\mathcal F'$, and its number of values is at most $|\sigma(H)|\cdot|\langle\mathcal F'\rangle(H)|$.
\end{proposition}

\begin{proof}
The condition $\mathrm{Serv}(\sigma)\subseteq\mathrm{Serv}(\tau)$ holds if and only if $\ker\tau\subseteq\ker\sigma$: the implication from right to left is immediate, and for the converse take the task $\sigma$ itself, which lies in $\mathrm{Serv}(\sigma)$ and so in $\mathrm{Serv}(\tau)$, giving $\ker\tau\subseteq\ker\sigma$. Together with $\ker\tau\subseteq\ker\mathcal F'$ this gives $\ker\tau\subseteq\ker\sigma\cap\ker\mathcal F'$, and the quotient by this intersection satisfies both conditions. The classes of $\ker\sigma\cap\ker\mathcal F'$ are the nonempty intersections of a class of $\ker\sigma$ with a class of $\ker\mathcal F'$, whence the bound.
\end{proof}

The extension is the pairing $\langle\sigma,\langle\mathcal F'\rangle\rangle$ restricted to its image. Nothing cheaper is available, and the bound is generally not attained: whenever the new tasks are partly determined by $\sigma$ the number of classes is smaller. When the new family is a function of old data, $\ker\sigma\subseteq\ker\mathcal F'$, the extension has no cost.

\section{Lower bounds}

\begin{theorem}[Cardinality bound]
\label{thm:lower}
If $\sigma$ is $\mathcal F$-sufficient then $|\sigma(H)|\ge|H/\ker\mathcal F|=|\langle\mathcal F\rangle(H)|$. In particular, any binary encoding of an $\mathcal F$-sufficient representation of a finite $H$ has length at least $\lceil\log_2|\langle\mathcal F\rangle(H)|\rceil$ for some history, and this is attained by $q_{\mathcal F}$.
\end{theorem}

\begin{proof}
By Theorem~\ref{thm:suff-refine}(ii) there is a surjection $u:\sigma(H)\to H/\ker\mathcal F$, so the domain has at least as many elements. The quotient $H/\ker\mathcal F$ is in bijection with the image of $\langle\mathcal F\rangle$. A binary code assigning distinct strings to $N$ values has a string of length at least $\lceil\log_2 N\rceil$ for some value, by the pigeonhole principle; the quotient attains $\lceil\log_2 N\rceil$ for fixed-length codes.
\end{proof}

For Example~\ref{ex:deriv-rep} the bound gives $2^k$ classes and $k$ bits, whatever else is stored. The derivation, which contributes zero bits to this family, is entirely irrelevant to it.

The same bound holds in expectation under a probability distribution on $H$, which is the natural setting if the family is to be served for an unknown, randomly arising run. Distributions are not given in this book and none is invented; the statement below holds for every distribution and is used only qualitatively.

\begin{theorem}[Entropic characterization]
\label{thm:entropy}
Let $P$ be a probability distribution on a finite set $H$, let $S=\sigma(h)$ and $T=f(h)$ for $h\sim P$. Then $f$ is $\sigma$-computable $P$-almost surely, that is, there is $f'$ with $f'(S)=T$ almost surely, if and only if $H(T\mid S)=0$. In that case $H(S)\ge H(T)$, with equality if and only if $S$ is also a function of $T$ almost surely.
\end{theorem}

\begin{proof}
If $T=f'(S)$ almost surely then $T$ is determined by $S$ and $H(T\mid S)=0$. If $H(T\mid S)=0$ then for each $s$ with positive probability the conditional distribution of $T$ is a point mass, and $f'(s)$ may be taken as its support; elsewhere $f'$ is arbitrary. By the chain rule
\[
H(S)=H(S,T)-H(T\mid S)=H(S,T)=H(T)+H(S\mid T)\ \ge\ H(T),
\]
with equality if and only if $H(S\mid T)=0$.
\end{proof}

The information-theoretic version of the lossy case, which does not appear in the remainder of the book except as a remark, is the data-processing inequality $I(T;g(S))\le I(T;S)$ for any function $g$ \cited{Cover and Thomas, \emph{Elements of Information Theory}, Theorem~2.8.1}. It shows that sufficiency, in the sense $I(T;S)=H(T)$, cannot be restored by further processing once lost.

\section{No free sufficiency}

The preceding results are meaningful only if some tasks are not served by compressed representations. The next theorem makes that unconditional.

\begin{theorem}[No free sufficiency]
\label{thm:nofree}
Let $\sigma:H\to R$. The following are equivalent.
\begin{enumerate}[label=(\roman*),nosep]
\item $\sigma$ is injective.
\item $\sigma$ is $\mathcal F$-sufficient for the family $\mathcal F_{\mathrm{all}}$ of all functions $H\to\{0,1\}$.
\item $\sigma$ is $\mathcal F$-sufficient for the family $\mathcal F_{\mathrm{pt}}=\{\mathbf 1_{\{h\}}:h\in H\}$ of point indicators.
\end{enumerate}
In particular, if $\sigma$ identifies two distinct histories then there is a task, indeed a $\{0,1\}$-valued one, that $\sigma$ does not serve.
\end{theorem}

\begin{proof}
(i)$\Rightarrow$(ii): if $\sigma$ is injective, any $f$ factors as $f\circ\sigma^{-1}$ on the image and arbitrarily elsewhere. (ii)$\Rightarrow$(iii) is trivial since $\mathcal F_{\mathrm{pt}}\subseteq\mathcal F_{\mathrm{all}}$. (iii)$\Rightarrow$(i): suppose $\sigma(h)=\sigma(h')$ with $h\neq h'$. Then $\mathbf 1_{\{h\}}(h)=1\neq0=\mathbf 1_{\{h\}}(h')$, so $\ker\sigma\not\subseteq\ker\mathbf 1_{\{h\}}$.
\end{proof}

\begin{corollary}[Compression is a bet on a family]
\label{cor:bet}
Every non-injective representation of the runs of a search procedure is sufficient for some declared families and insufficient for others. The set of families it is sufficient for is exactly $\{\mathcal F:\ker\sigma\subseteq\ker\mathcal F\}$, an up-closed set in the lattice of families under the order by kernel inclusion.
\end{corollary}

The corollary is the formal justification for requiring a declaration. An artifact that carries only a derivation is a bet that no later task will distinguish runs with the same derivation. The bet is often correct for checking and for citation; it is not correct for tasks whose answer depends on the discarded search record, such as asking which strengthening of a hypothesis would have been cheaper to prove, or why a related goal did not succeed under the same strategy.

\section{Proof irrelevance as a sufficiency statement}

Chapter~\ref{ch:kernels} noted that proof irrelevance in a type theory with a proof-irrelevant sort of propositions collapses all proofs of the same proposition. The language of this chapter turns that observation into a theorem about which tasks survive.

\begin{definition}[Proof-irrelevant representation]
\label{def:pirrel}
Let $H$ be a set of pairs $(P,p)$ of a proposition $P$ and a proof term $p$ of $P$ (up to definitional equality). The \emph{proof-irrelevant representation} is $\pi(P,p)=P$.
\end{definition}

\begin{proposition}[Tasks served by proof irrelevance]
\label{prop:pirrel}
A task $f$ is $\pi$-sufficient if and only if $f(P,p)=f(P,p')$ for all $(P,p),(P,p')\in H$, that is, if and only if its answer depends on the proposition alone. In particular:
\begin{enumerate}[label=(\roman*),nosep]
\item type checking of any term that mentions the proof only through its type is $\pi$-sufficient;
\item the task ``return the algorithm encoded by $p$'' is not $\pi$-sufficient whenever two proofs of some $P$ encode different algorithms;
\item the task ``does $p$ use the axiom $a$'' is not $\pi$-sufficient in general, so axiom footprints are not an invariant of the proposition.
\end{enumerate}
\end{proposition}

\begin{proof}
The first statement is Proposition~\ref{prop:factor} for $\pi$. For (i), a term of the form $c\,P\,p$ is typed by the type of $p$ alone, by the definition of the typing rule for a proof-irrelevant argument. For (ii) and (iii), let $P$ have proofs $p,p'$ with different footprints; Chapter~\ref{ch:kernels} gives such pairs, for instance a proof that uses classical choice and one that does not. The task is non-constant on the fibre over $P$.
\end{proof}

This is the precise sense in which proof irrelevance is an \emph{admissible collapse} for dependents that consume a theorem by its statement and an \emph{inadmissible collapse} for any procedure that must recover what was done. It is not that one is correct and the other wrong; they are tasks with different kernels, and the declared family decides.

\section{A geometric instance: exterior data and hidden material}
\label{sec:geometric}

The abstract statements of this chapter have a concrete instance in which the extraction map, its fibres, and a decidable certificate for the nonemptiness of a fibre can all be written down. A flat sheet is mapped into space; the extraction keeps only what the map does on a marked part of the sheet; the question is what that restriction fixes about the rest.

\begin{definition}[Marking problem]
\label{def:marking}
Let $\Sigma\subset\mathbb R^2$ be a compact convex set with nonempty interior, the \emph{sheet}. A \emph{marking} is a compact set $A=A_1\cup\dots\cup A_k\subset\Sigma$ with the $A_i$ pairwise disjoint, together with a \emph{target} $\tau:A\to\mathbb R^3$ whose restriction to each $A_i$ preserves distances. A \emph{path isometry} is a continuous map $f:\Sigma\to\mathbb R^3$ such that $f\circ\gamma$ has the same length as $\gamma$ for every rectifiable path $\gamma$ in $\Sigma$; every $C^1$ isometric immersion is one. For $\delta\ge0$, a \emph{$\delta$-realization} of $\tau$ is a path isometry $f$ with $|f(p)-\tau(p)|\le\delta$ for $p\in A$; a $0$-realization is a \emph{realization}. Write $\mathcal R_\delta(\tau)$ for the set of $\delta$-realizations. The \emph{exterior map} is $E(f)=f|_A$, so that $\mathcal R_0(\tau)=E^{-1}(\tau)$ is a fibre of $E$ on the set of path isometries.
\end{definition}

Path isometries may self-intersect. Nothing below concerns embeddedness, which is a stricter question.

\begin{definition}[Nonexpansion condition]
\label{def:L}
The target $\tau$ satisfies $(\mathrm L)$ if $|\tau(p)-\tau(q)|\le|p-q|$ for all $p,q\in A$, and $(\mathrm L^{<})$ if the inequality is strict whenever $p,q$ lie in different $A_i$. Since $\tau$ preserves distances on each $A_i$, equality holds within each square, so $(\mathrm L)$ is a condition on cross pairs.
\end{definition}

\begin{lemma}[Necessity]
\label{lem:geo-nec}
If $f$ is a $\delta$-realization of $\tau$, then $|\tau(p)-\tau(q)|\le|p-q|+2\delta$ for $p,q\in A$. In particular $\mathcal R_0(\tau)\neq\emptyset$ implies $(\mathrm L)$.
\end{lemma}

\begin{proof}
The segment from $p$ to $q$ lies in $\Sigma$ by convexity and has length $|p-q|$. Its image under $f$ has the same length and joins $f(p)$ to $f(q)$, so $|f(p)-f(q)|\le|p-q|$. Then $|\tau(p)-\tau(q)|\le|f(p)-f(q)|+2\delta$.
\end{proof}

\begin{theorem}[Realizability criterion]
\label{thm:geo-criterion}
The following are equivalent: (a) $\mathcal R_\delta(\tau)\neq\emptyset$ for every $\delta>0$; (b) $(\mathrm L)$ holds. Moreover, when they hold, the $\delta$-realizations may be chosen with image in the $\delta$-neighbourhood of the convex hull $\mathrm{conv}\,\tau(A)$.
\end{theorem}

\begin{proof}
(a)$\Rightarrow$(b) is Lemma~\ref{lem:geo-nec} with $\delta\to0$. For (b)$\Rightarrow$(a), the hypothesis makes $\tau$ a $1$-Lipschitz map from $A\subset\mathbb R^2$ to $\mathbb R^3$: within a square it preserves distances, and across squares it is nonexpanding by $(\mathrm L)$. By Kirszbraun's theorem (Appendix~\ref{app:external}, A.23) $\tau$ extends to a $1$-Lipschitz map $F:\mathbb R^2\to\mathbb R^3$. Let $K=\mathrm{conv}\,\tau(A)$ and $\pi_K$ the nearest-point projection to $K$, which is $1$-Lipschitz and fixes $\tau(A)$; replace $F$ by $\pi_K\circ F$, so $F(\mathbb R^2)\subset K$. Fix $c\in K$ and $s\in(0,1)$ and put $G=(1-s)F+sc$, a $(1-s)$-Lipschitz map into $K$ with $|G-F|\le s\,\mathrm{diam}\,K$. Mollify $G$ by a kernel of radius $\eta$: the result $G_\eta$ is smooth, still $(1-s)$-Lipschitz, takes values in $K$ (mollification averages points of $K$), and $|G_\eta-G|\le(1-s)\eta$ on $\Sigma$. Its differential satisfies $\|dG_\eta\|\le1-s<1$, so $G_\eta$ is a strictly short smooth map on a smooth compact convex neighbourhood $\tilde\Sigma\supset\Sigma$. By the Nash--Kuiper theorem in Gromov's form (A.24), for every $\rho>0$ there is a $C^1$ isometric immersion $f:\tilde\Sigma\to\mathbb R^3$ with $|f-G_\eta|<\rho$. Restricted to $\Sigma$, $f$ is a path isometry, and on $A$
\[
|f-\tau|\le\rho+\eta+s\,\mathrm{diam}\,K\le\delta
\]
for $\rho,\eta,s$ chosen small in terms of $\delta$. The image of $f$ lies within $\rho$ of $K$.
\end{proof}

Exact attainment on $A$ is a different matter, since it asks for a path isometry that agrees with $\tau$ on $A$ and is strictly short nowhere that matters near $\partial A$; it needs a relative form of the isometric immersion theorem and is not claimed here.

The criterion is a \emph{certificate of satisfiability} in the sense of Chapter~\ref{ch:certificates}: the target $\tau$ together with the verification of $(\mathrm L)$ certifies that a hidden extension exists, and does so without exhibiting it.

\begin{proposition}[Decidability]
\label{prop:geo-decid}
If the $A_i$ are polygons and $\tau|_{A_i}$ are affine, then $(\mathrm L)$ is a universal sentence of real closed fields and is decidable; the three-valued verdict \textsc{realizable} ($\mathrm L$ holds, Theorem~\ref{thm:geo-criterion}), \textsc{not realizable} ($\mathrm L$ fails, Lemma~\ref{lem:geo-nec}) has no third value.
\end{proposition}

\begin{proof}
Each $A_i$ is a finite union of sets defined by linear inequalities, $\tau$ is given by rational affine maps, and $(\mathrm L)$ is $\forall p\,\forall q\,(p\in A\wedge q\in A\to|\tau p-\tau q|^2\le|p-q|^2)$, a first-order sentence of real arithmetic. Tarski's theorem (A.25) decides it. If $(\mathrm L)$ fails, then $|\tau(p)-\tau(q)|>|p-q|$ for some pair, and $2\delta$ below the excess gives a contradiction with Lemma~\ref{lem:geo-nec}, so $\mathcal R_\delta(\tau)=\emptyset$ for small $\delta$.
\end{proof}

In the vocabulary of Chapter~\ref{ch:inheritance} this is a certified fragment in which preservation is decidable and the verdict \textsc{outside the certified fragment} is never needed.

\begin{example}[Three faces at a corner]
\label{ex:corner}
Let $g>0$ and take three unit squares in the plane,
\[
A_1=\{(-z-g,\,y)\},\quad A_2=\{(x,\,-z-g)\},\quad A_3=[0,1]^2,\qquad x,y,z\in[0,1],
\]
mapped to three faces of the unit cube meeting at the origin by $\tau(-z-g,y)=(0,y,z)$, $\tau(x,-z-g)=(x,0,z)$, $\tau(x,y)=(x,y,0)$. Each restriction preserves distances. For $p=(-z-g,y_1)\in A_1$ and $q=(x',-z'-g)\in A_2$ expanding the squares gives
\[
|p-q|^2-|\tau p-\tau q|^2=2x'(z+g)+2y_1(z'+g)+2zz'+2g(z+z')+2g^2\ \ge\ 2g^2,
\]
and for $p\in A_1$, $q=(x,y)\in A_3$,
\[
|p-q|^2-|\tau p-\tau q|^2=2xz+2g(x+z)+g^2\ \ge\ g^2,
\]
with the same bound for $A_2$ and $A_3$ by symmetry. Hence $(\mathrm L^{<})$ holds with slack at least $g^2$ in squared distance, and by Theorem~\ref{thm:geo-criterion} the three faces can be exposed, up to any tolerance, by a path isometry of any sheet containing them, with the material between and around them accommodated in the convex hull of the three faces. The squares need not be adjacent in the sheet; the gap $g$ is the price of their separation.
\end{example}

The remaining question is what the exterior data leave undetermined.

\begin{theorem}[The exterior does not determine the interior]
\label{thm:geo-residue}
Assume $(\mathrm L)$ and let $p_0\in\Sigma\setminus A$. Put
\[
\Phi(p_0)=\{y\in\mathbb R^3:\ |y-\tau(q)|<|p_0-q|\ \text{for all } q\in A\}.
\]
Then $\Phi(p_0)$ is open and convex; for every $y\in\Phi(p_0)$ and $\delta>0$ there is a $\delta$-realization $f$ with $|f(p_0)-y|\le\delta$; and every $\delta$-realization satisfies $|f(p_0)-\tau(q)|\le|p_0-q|+2\delta$ for all $q\in A$, so that an exact realization has $f(p_0)$ in the closed set $\bigcap_{q}\overline B(\tau(q),|p_0-q|)\supseteq\Phi(p_0)$. In particular, if $\Phi(p_0)\neq\emptyset$ then the value $f(p_0)$ is not a function of the exterior data: there are $\delta$-realizations with the same exterior map up to $\delta$ and values at $p_0$ at arbitrary distance up to $\mathrm{diam}\,\Phi(p_0)$.
\end{theorem}

\begin{proof}
Convexity holds because each set $\{y:|y-\tau(q)|<r\}$ is a convex ball. Openness: $y\mapsto\max_{q\in A}(|y-\tau(q)|-|p_0-q|)$ is a continuous function on $\mathbb R^3$ (the maximum of a continuous function over the compact $A$), and $\Phi(p_0)$ is where it is negative. For the second statement, set $A'=A\cup\{p_0\}$ with $A_{k+1}=\{p_0\}$ and $\tau'(p_0)=y$. The restriction to a single point preserves distances, and $(\mathrm L)$ holds for $\tau'$: pairs in $A$ are covered by the hypothesis and pairs $(p_0,q)$ by $y\in\Phi(p_0)$. Theorem~\ref{thm:geo-criterion} gives $\delta$-realizations of $\tau'$, which are $\delta$-realizations of $\tau$ with $|f(p_0)-y|\le\delta$. The inequality $|f(p_0)-\tau(q)|\le|p_0-q|+2\delta$ is Lemma~\ref{lem:geo-nec} applied to the pair $(p_0,q)$ for a $\delta$-realization, using the path-isometry bound $|f(p_0)-f(q)|\le|p_0-q|$. For the non-function statement, pick $y_1\ne y_2$ in $\Phi(p_0)$ and $\delta<|y_1-y_2|/2$; the realizations $f_1,f_2$ so obtained have the same exterior up to $\delta$ and $f_1(p_0)\ne f_2(p_0)$. If $f(p_0)$ were a function of the exterior data it would be constant on the fibre, contradicting Proposition~\ref{prop:factor}.
\end{proof}

\begin{example}
\label{ex:corner-residue}
In Example~\ref{ex:corner} take $p_0=(-g,-g)$, which lies in the gap between the three squares, and $y=(0,0,0)$. For $q=(-z-g,y_1)\in A_1$ one has $|p_0-q|^2=z^2+(g+y_1)^2>z^2+y_1^2=|y-\tau(q)|^2$; for $A_2$ the same by symmetry; for $q=(x,y_1)\in A_3$, $|p_0-q|^2=(g+x)^2+(g+y_1)^2>x^2+y_1^2$. So $y\in\Phi(p_0)$, and $\Phi(p_0)$ is a nonempty open set. The exterior data, the three exposed faces, leave the position of the sheet's gap point free in an open region of space.
\end{example}

By Theorem~\ref{thm:geo-residue} the residue of the exterior map has positive dimension whenever there is slack, and the slack $\min_q(|p_0-q|-|y-\tau(q)|)$ is its quantitative size. That is the same quantity as the margin of a Farkas certificate in Chapter~\ref{ch:history-repair}: a certified region has room, and the room is what the certificate does not pin down. Conversely, the closed region $C(p_0)=\bigcap_q\overline B(\tau(q),|p_0-q|)$ is exactly the set of positions attainable at $p_0$, and by strict convexity of the balls it is either a single forced point or has nonempty interior; intermediate dimensions do not occur. (For two distinct points of $C(p_0)$, the midpoint satisfies $|m-\tau(q)|<|p_0-q|$ for every $q$, since the Euclidean norm is strictly convex and two distinct points at equal distance from $\tau(q)$ cannot both lie on one ray from it; so the midpoint lies in the open set $\Phi(p_0)$.) Wherever $C(p_0)$ has more than one point, a task that depends on $f(p_0)$, such as locating a crease or deciding which layer lies above which, is therefore not served by the exterior map, by Proposition~\ref{prop:factor}, and is served by any extraction that retains $f(p_0)$.

\section{Declaring families in practice}

A declaration should name tasks, not representations. Table~\ref{tab:families} lists task families used in the remainder of the book, with the kernel each induces on a run of a proof search and the smallest standard artifact that is sufficient for it.

\begin{table}[ht]
\centering
\small
\begin{tabular}{@{}p{3.4cm}p{4.6cm}p{4.2cm}@{}}
\toprule
Family & Distinction it detects & Smallest standard sufficient artifact\\
\midrule
Check & Derivation validity & Checked derivation\\
Cite & Statement and axiom footprint & Statement and footprint\\
Replay & Derivation under changed implementation & Derivation plus kernel version\\
Adapt to a nearby goal & Steps that depend on the goal's details & Derivation with dependency marking (Ch.~\ref{ch:history-repair})\\
Diagnose a failure & Which options were refuted, and by what & Run trace restricted to refusals (Ch.~\ref{ch:search-worlds})\\
Audit scope & Which assumptions entered the correspondence claim & Obligation ledger (Ch.~\ref{ch:obligations})\\
\bottomrule
\end{tabular}
\caption{Task families and the minimal artifacts that serve them. The last column states what is sufficient, not what is conventionally published.}
\label{tab:families}
\end{table}

The table's entries are definitions of family names, not empirical claims, with one exception. The assertion that the conventional publication of a result, a statement together with a derivation, is not sufficient for the diagnosis and audit families is a consequence of Theorem~\ref{thm:nofree} and a nonconstancy check on a single example, as in Example~\ref{ex:deriv-rep}. Chapters~\ref{ch:history-repair} and~\ref{ch:mutation} supply the examples for the remaining rows.

\section{What this chapter fixes}

\begin{principal}
\textbf{Sufficiency is always sufficiency for something.} (1) A representation serves a task if and only if the representation's kernel refines the task's kernel (Theorem~\ref{thm:suff-refine}). (2) Declaring a family fixes a minimal artifact and a closure of served tasks (Theorem~\ref{thm:closure}). (3) The cost of any sufficient representation is bounded below by the number of classes of the family kernel (Theorem~\ref{thm:lower}). (4) No non-injective representation serves all tasks (Theorem~\ref{thm:nofree}).
\end{principal}

The chapters that follow ask how a representation fares under \emph{change}, which requires more than a family of tasks over a fixed history space: it requires maps between history spaces, and a notion of which of them a representation survives. That is the subject of Chapter~\ref{ch:history-repair}.

\begin{exercise}
Show that the set of $\mathcal F$-sufficient kernels is closed under common refinement but, in general, not under the join in $\Pi(H)$. Give a four-element example.
\end{exercise}

\begin{exercise}
Prove that $\mathrm{Serv}(\mathrm{Rep}(\mathrm{Serv}(\sigma)))=\mathrm{Serv}(\sigma)$ and that $\mathrm{Rep}(\mathrm{Serv}(\mathrm{Rep}(\mathcal F)))=\mathrm{Rep}(\mathcal F)$.
\end{exercise}

\begin{exercise}
In Example~\ref{ex:deriv-rep} let each $A$ occur with exactly one $d$. Compute $\ker\sigma\cap\ker\mathcal F_{\mathrm{ref}}$ and state the cost of the cheapest extension of $\sigma$ to $\mathcal F_{\mathrm{ref}}$ using Proposition~\ref{prop:extend}.
\end{exercise}
